10912: \subsection{Unbounded source ratio and relative current}
10913: \label{component:pef021eabacde}
10914: 
10915: \label{pef021eabacde:sec:r2-unbounded-current-new}
10916: 
10917: \paragraph{Precise assertion.}
10918: Retain the original RER source, all its strict gaps, and the notation
10919: TU2--TU16 of the preceding bounded-relative-current proof.
10920: Assume the native $a,t$ are positive unbounded and
10921: \[
10922:              \lambda=\frac{n}{tB}=\frac{n}{qz}
10923:                    \text{ is positive unbounded}.                 \tag{UL1}
10924: \]
10925: Then the whole original weak-C2 parent is tensor SOS for every
10926: independent projective test and every state. The argument below
10927: complements TU1; it never assigns a sign to the formal C1 coefficient
10928: $\beta^0$, and never replaces the original physical source by that
10929: formal source.
10930: 
10931: \subsubsection*{Source bounds without a thinness assumption}
10932: 
10933: The exact RER physicality equation already used in TU4 is
10934: \[
10935:  Q_{\rm RER}=\tau\{A+2r(z\Delta)\lambda\},\qquad
10936:  \tau=\frac q{zw},\qquad
10937:  A=(1+z)(1+Sh)(1+Sq).
10938: \]
10939: Here $A,Q_{\rm RER},r,z\Delta$ are positive units with residue one.
10940: Since $\lambda$ is unbounded, this proves, with no cancellation of
10941: mixed factors,
10942: \[
10943:  \tau\in\mathrm I,\qquad \lambda\tau\sim\frac12,\qquad
10944:  \lambda q\sim\frac{zw}{2},\qquad
10945:  \frac n{z^2w}\sim\frac12.                                     \tag{UL2}
10946: \]
10947: All four divisions are pure source. Recall $0<w<R\sim1$.
10948: The exact TU9 formulas therefore imply
10949: \[
10950: \begin{gathered}
10951:  (1-h)\delta c_b,\ h\delta a_b,\ \delta D_b,\
10952:             \delta h_2,\delta h_b=\lambda q\,O(1),\\
10953:  \delta\alpha_b/q^2=\lambda q\,O(S+zw),\qquad
10954:  h_2\asymp zw,\quad h_b=zw/2,\\
10955:  H_0=4h_2/R+2z\asymp z,\qquad H_0\geq2z,\\
10956:  \frac g{\sqrt{H_0}}=O(1),\quad
10957:  \frac c{\sqrt{H_0}}=O(1),\quad
10958:  \frac{\lambda q}{H_0}=O(1),\quad \frac n{z^2}=O(1).
10959: \end{gathered}                                                   \tag{UL3}
10960: \]
10961: In particular $\lambda q$ is infinitesimal. The displayed coefficient
10962: arrays have denominators only among the pure-source units
10963: $P,J,Q,R$ after the original $h,1-h$ openings cancel.
10964: The actual and formal $\sigma=c_b+a_b$ differ by a relative
10965: infinitesimal. Their square-root ratio is a bounded unit.
10966: No lower bound on $w,h,1-h$ or the strict gap $M$ is assumed.
10967: 
10968: Use the complete generic SC3 identity with the actual coefficients.
10969: Retain its positive $m$ positive square term and use $\widehat N=\sqrt{H_0}N$.
10970: The formal comparison is used only for finite scalar and normal
10971: identities. Its $q$-coquotients and the QI coefficient arrays are
10972: bounded here: their derivation requires $S,z,q$ infinitesimal and
10973: $w$ bounded, but not $\beta^0\geq0$. The sole QI bound that is no
10974: longer available, $w/(c/g)=O(1)$, is replaced explicitly below.
10975: 
10976: \subsubsection*{The entire extra current supplies a new angular pivot}
10977: 
10978: Use the exact positive-chart rows TU15--TU16, denoted here by
10979: $A_*,B_*,C_*,D_*,U_*$ and $\Gamma$. Define
10980: \[
10981:  J_3=\sqrt\lambda(U_*-B_*),\qquad
10982:  E_*=h_t-v,\qquad F_*=vh_t-gD_*.
10983:                                                                \tag{UL4}
10984: \]
10985: The first line of the exact whole-current identity TU17 is
